\documentclass[10pt,a4paper]{article}
\usepackage[utf8]{inputenc}
\usepackage[T1]{fontenc}
\usepackage[ngerman]{babel}
\usepackage{helvet}
\renewcommand{\familydefault}{\sfdefault}
\usepackage[left=1cm, right=1cm, top=1cm, bottom=1cm]{geometry}
\usepackage{amsmath,amssymb}
\usepackage{array}
\usepackage{colortbl}
\usepackage{multicol}
\usepackage{tcolorbox}
\tcbuselibrary{skins}

% Farben
\definecolor{physikblue}{RGB}{44,90,160}
\definecolor{lightblue}{RGB}{220,235,250}
\definecolor{merkegreen}{RGB}{34,139,94}

\pagestyle{empty}
\setlength{\parindent}{0pt}

\begin{document}

\begin{center}
{\Large\bfseries\color{physikblue} Formelkarte Kinematik} \hfill {\small Klasse 10E}
\end{center}

\vspace{-0.3cm}
\hrule height 1pt
\vspace{0.3cm}

\begin{multicols}{2}

% === GLEICHFÖRMIGE BEWEGUNG ===
\begin{tcolorbox}[colback=lightblue, colframe=physikblue, boxrule=0.5pt,
    title={\small\bfseries Gleichförmige Bewegung}, left=2pt, right=2pt, top=2pt, bottom=2pt]
\small
\[
v = \frac{s}{t} \quad\Rightarrow\quad s = v \cdot t \quad\Rightarrow\quad t = \frac{s}{v}
\]
Mit Startposition: $s(t) = v \cdot t + s_0$
\end{tcolorbox}

\vspace{0.2cm}

% === BESCHLEUNIGTE BEWEGUNG ===
\begin{tcolorbox}[colback=lightblue, colframe=physikblue, boxrule=0.5pt,
    title={\small\bfseries Beschleunigte Bewegung}, left=2pt, right=2pt, top=2pt, bottom=2pt]
\small
\textbf{Definition:} $a = \dfrac{\Delta v}{\Delta t} = \dfrac{v - v_0}{t}$

\vspace{0.2cm}
\textbf{Grundformeln:}
\begin{align*}
v(t) &= a \cdot t + v_0 \\
s(t) &= \frac{1}{2} a \cdot t^2 + v_0 \cdot t + s_0
\end{align*}

\textbf{Zeitlose Formel:}
\[
v^2 = v_0^2 + 2 \cdot a \cdot s
\]

\textbf{Aus Ruhe} ($v_0 = 0$, $s_0 = 0$):
\[
v = a \cdot t \qquad s = \frac{1}{2} a \cdot t^2
\]
\end{tcolorbox}

\vspace{0.2cm}

% === BREMSEN ===
\begin{tcolorbox}[colback=lightblue, colframe=physikblue, boxrule=0.5pt,
    title={\small\bfseries Bremsvorgang}, left=2pt, right=2pt, top=2pt, bottom=2pt]
\small
Bei Verzögerung: $a < 0$

\[
t_{\text{brems}} = \frac{v_0}{|a|} \qquad s_{\text{brems}} = \frac{v_0^2}{2|a|}
\]

\textit{Merke: Doppelte Geschwindigkeit = 4-facher Bremsweg!}
\end{tcolorbox}

\columnbreak

% === FREIER FALL ===
\begin{tcolorbox}[colback=merkegreen!15, colframe=merkegreen, boxrule=0.5pt,
    title={\small\bfseries Freier Fall}, left=2pt, right=2pt, top=2pt, bottom=2pt]
\small
\textbf{Fallbeschleunigung:} $g \approx 10\,\mathrm{m/s^2}$

\vspace{0.1cm}
\begin{tabular}{ll}
Geschwindigkeit: & $v = g \cdot t = 10t$ \\
Fallstrecke: & $s = \frac{1}{2}g \cdot t^2 = 5t^2$ \\
Fallzeit: & $t = \sqrt{\frac{s}{5}}$ \\
Aufprall: & $v = \sqrt{2gs}$
\end{tabular}

\vspace{0.2cm}
\textbf{Senkrechter Wurf nach oben:}\\
$v(t) = v_0 - g \cdot t$ \quad (bremst!)\\
Steigzeit: $t_{\text{steig}} = \frac{v_0}{g}$\\
Steighöhe: $h_{\text{max}} = \frac{v_0^2}{2g}$
\end{tcolorbox}

\vspace{0.2cm}

% === UMRECHNUNGEN ===
\begin{tcolorbox}[colback=lightblue, colframe=physikblue, boxrule=0.5pt,
    title={\small\bfseries Umrechnungen}, left=2pt, right=2pt, top=2pt, bottom=2pt]
\small
\textbf{km/h $\leftrightarrow$ m/s:}
\begin{center}
\begin{tabular}{|c|c|c|c|c|}
\hline
\textbf{km/h} & 36 & 72 & 90 & 108 \\
\hline
\textbf{m/s} & 10 & 20 & 25 & 30 \\
\hline
\end{tabular}
\end{center}
km/h $\rightarrow$ m/s: $\div 3{,}6$ \quad m/s $\rightarrow$ km/h: $\times 3{,}6$

\vspace{0.2cm}
\textbf{Kopfrechenbare Wurzeln:}\\
$\sqrt{4}=2$, $\sqrt{9}=3$, $\sqrt{16}=4$, $\sqrt{25}=5$,\\
$\sqrt{36}=6$, $\sqrt{49}=7$, $\sqrt{64}=8$, $\sqrt{81}=9$
\end{tcolorbox}

\vspace{0.2cm}

% === DIAGRAMME ===
\begin{tcolorbox}[colback=lightblue, colframe=physikblue, boxrule=0.5pt,
    title={\small\bfseries Diagramme lesen}, left=2pt, right=2pt, top=2pt, bottom=2pt]
\small
\begin{tabular}{|l|c|c|}
\hline
\textbf{Diagramm} & \textbf{Steigung} & \textbf{Fläche} \\
\hline
$s(t)$ & $v$ & --- \\
\hline
$v(t)$ & $a$ & $s$ (Weg!) \\
\hline
$a(t)$ & --- & $\Delta v$ \\
\hline
\end{tabular}
\end{tcolorbox}

\end{multicols}

\vspace{0.2cm}
\hrule height 0.5pt
\vspace{0.2cm}

% === ANALOGIE-ÜBERSICHT ===
\begin{center}
\small
\renewcommand{\arraystretch}{1.2}
\begin{tabular}{|l|c|c|c|c|}
\hline
\rowcolor{lightblue}
\textbf{Bewegungsart} & \textbf{Beschleunigung} & \textbf{Geschwindigkeit} & \textbf{Weg} & \textbf{Zeitlos} \\
\hline
Gleichförmig & $a = 0$ & $v = \text{konst.}$ & $s = v \cdot t$ & --- \\
\hline
Beschleunigt (aus Ruhe) & $a = \text{konst.}$ & $v = a \cdot t$ & $s = \frac{1}{2}at^2$ & $v^2 = 2as$ \\
\hline
Freier Fall & $a = g \approx 10$ & $v = g \cdot t$ & $s = 5t^2$ & $v^2 = 2gs$ \\
\hline
Wurf nach oben & $a = -g$ & $v = v_0 - gt$ & $s = v_0 t - 5t^2$ & $v^2 = v_0^2 - 2gs$ \\
\hline
\end{tabular}
\end{center}

\vspace{0.3cm}

\begin{center}
\fbox{\parbox{0.9\textwidth}{\centering\small
\textbf{Luftwiderstand (nur qualitativ):} $F_W \sim v^2$ \quad
Größere Fläche $\Rightarrow$ mehr Widerstand $\Rightarrow$ kleinere $v_{\text{grenz}}$\\
Grenzgeschwindigkeit erreicht bei $F_G = F_W$
}}
\end{center}

\end{document}
